\section{光滑函数的稠密性}
\label{sec:sobolev-smooth-density}

Sobolev函数可以有折点，也可以无界，直接对它们做计算往往不方便。常用的办法是先对光滑函数证明某个等式或估计，再用光滑函数逼近一般的Sobolev函数，让结论在极限中保留下来。这要求光滑函数在Sobolev范数下稠密：不仅函数本身要接近，每个弱导数分量也要同时接近。本节对任意开集和有限指数证明这一点，所用工具是第~\ref{chap:smoothing-weak-derivatives}章的磨光与第~\ref{chap:function-approximation}章的$L^p$磨光收敛。

\paragraph{需要的两条磨光结论}
记$\rho_\varepsilon(x)=\varepsilon^{-d}\rho(x/\varepsilon)$为第~\ref{chap:smoothing-weak-derivatives}章的磨光核：$\rho\geq0$光滑，支撑在单位闭球内，积分为$1$。对局部可积的$f$，记$\Omega_\varepsilon=\{x\in\Omega:\operatorname{dist}(x,\partial\Omega)>\varepsilon\}$，磨光$f_\varepsilon=f*\rho_\varepsilon$在$\Omega_\varepsilon$上光滑。若$f$有局部可积的弱导数$\partial^\alpha f$，则在$\Omega_\varepsilon$上
\begin{equation}
\partial^\alpha(f*\rho_\varepsilon)=(\partial^\alpha f)*\rho_\varepsilon ,
\label{eq:sobolev-mollifier-commute}
\end{equation}
左端是光滑函数的经典导数，右端是弱导数的磨光。这是第~\ref{chap:smoothing-weak-derivatives}章的结论。另一条结论来自第~\ref{chap:function-approximation}章：在全空间上，对$1\leq p<\infty$和$g\in L^p(\mathbb R^d)$，有$\lVert g*\rho_\varepsilon-g\rVert_p\to0$。

把两条结论用于$\Omega=\mathbb R^d$上的$f\in W^{k,p}(\mathbb R^d)$，此时$\Omega_\varepsilon=\mathbb R^d$。对每个$|\alpha|\leq k$，式\eqref{eq:sobolev-mollifier-commute}给出$\partial^\alpha f_\varepsilon-\partial^\alpha f=(\partial^\alpha f)*\rho_\varepsilon-\partial^\alpha f$，它在$L^p$中趋于零。有限多个分量都趋于零，所以
\begin{equation}
\lVert f*\rho_\varepsilon-f\rVert_{W^{k,p}(\mathbb R^d)}\to0\qquad(\varepsilon\to0,\ 1\leq p<\infty).
\label{eq:sobolev-whole-space-mollify}
\end{equation}
在一般开集上，磨光只在$\Omega_\varepsilon$上有定义，越靠近边界，允许的磨光半径越小。解决办法是把函数切成许多块，每块离边界的距离不同，各自选择磨光半径，再拼回去。切块需要光滑的截断函数与乘积法则。

\paragraph{截断函数}
先构造在给定紧集上取值为$1$、支撑仍在指定开集内的光滑函数。

\begin{lemma}[截断函数]
\label{lem:sobolev-cutoff}
设$K\subset V\subset\mathbb R^d$，$K$紧，$V$开。则存在$\eta\in C_c^\infty(V)$，使$0\leq\eta\leq1$，并且在$K$上$\eta=1$。
\end{lemma}

\begin{proof}
$V=\mathbb R^d$时取$\delta=1$；否则$K$到闭集$\mathbb R^d\setminus V$的距离为正，取$\delta$为这个距离的四分之一。令$U=\{x:\operatorname{dist}(x,K)<2\delta\}$，$\eta=\mathbf 1_U*\rho_\delta$。由磨光的光滑性，$\eta$光滑；由$\rho_\delta\geq0$且积分为$1$，$0\leq\eta\leq1$。若$x\in K$，球$B(x,\delta)\subset U$，于是$\eta(x)=\int_{B(x,\delta)}\rho_\delta(x-y)\,\mathrm dy=1$。若$\operatorname{dist}(x,K)>3\delta$，球$B(x,\delta)$与$U$不交，$\eta(x)=0$。所以$\eta$的支撑含于紧集$\{x:\operatorname{dist}(x,K)\leq3\delta\}\subset V$。
\end{proof}

称这样的$\eta$为\term{截断函数}（cutoff function）。用它乘一个函数，可以把函数限制在指定区域附近，而不引入不光滑的边缘。

\paragraph{乘积法则}
截断后的函数是否仍在Sobolev空间中，取决于光滑函数与Sobolev函数的乘积怎样求弱导数。

\begin{lemma}[光滑函数的乘积法则]
\label{lem:sobolev-product-rule}
设$f\in W^{k,p}(\Omega)$，$\zeta\in C^\infty(\Omega)$且$\zeta$及其各阶导数在$\Omega$上有界。则$\zeta f\in W^{k,p}(\Omega)$，一阶弱导数为$\partial_i(\zeta f)=\zeta\,\partial_if+(\partial_i\zeta)f$，并且存在只依赖$k,d$与$\zeta$的各阶导数上界的常数$C$，使$\lVert\zeta f\rVert_{W^{k,p}}\leq C\lVert f\rVert_{W^{k,p}}$。
\end{lemma}

\begin{proof}
对测试函数$\varphi$，$\zeta\varphi$仍是测试函数，经典乘积法则给出$\zeta\,\partial_i\varphi=\partial_i(\zeta\varphi)-(\partial_i\zeta)\varphi$。于是
\begin{equation*}
\begin{aligned}
\int_\Omega\zeta f\,\partial_i\varphi\,\mathrm dx&=\int_\Omega f\,\partial_i(\zeta\varphi)\,\mathrm dx-\int_\Omega f(\partial_i\zeta)\varphi\,\mathrm dx\\
&=-\int_\Omega\bigl(\zeta\,\partial_if+(\partial_i\zeta)f\bigr)\varphi\,\mathrm dx ,
\end{aligned}
\end{equation*}
第二个等号对测试函数$\zeta\varphi$用了$f$的弱导数定义。右端括号内的函数属于$L^p$，范数不超过$\sup|\zeta|\,\lVert\partial_if\rVert_p+\sup|\partial_i\zeta|\,\lVert f\rVert_p$。高阶情形对阶数归纳：$\zeta\,\partial_if$与$(\partial_i\zeta)f$都是“有界光滑函数乘以$W^{k-1,p}$函数”，对它们再用一阶结论。每一步只出现$\zeta$的有限个导数与$f$的有限个弱导数，范数界由此得到。
\end{proof}

\paragraph{局部有限的单位分解}
现在把$\Omega$按离边界的远近分层。令
\[
G_j=\{x\in\Omega:|x|<j,\ \operatorname{dist}(x,\mathbb R^d\setminus\Omega)>1/j\},\qquad j\geq1,
\]
$\Omega=\mathbb R^d$时省去距离条件，并约定$G_0=G_{-1}=\varnothing$。每个$G_j$是开集，闭包是$\Omega$内的紧集，$\overline{G_j}\subset G_{j+1}$，并且$\bigcup_jG_j=\Omega$。令$V_j=G_{j+1}\setminus\overline{G_{j-1}}$。这些开集覆盖$\Omega$，并且是\term{局部有限}（locally finite）的：$G_m$只与$V_1,\ldots,V_{m+1}$相交，所以每点有一个邻域只碰到有限多个$V_j$。

由引理~\ref{lem:sobolev-cutoff}，对紧集$\overline{G_j}\setminus G_{j-1}\subset V_j$取截断函数$\eta_j\in C_c^\infty(V_j)$，在该紧集上$\eta_j=1$。这些紧集覆盖$\Omega$，所以$S=\sum_i\eta_i$处处至少为$1$；由局部有限性，它在每点附近只是有限和，因而光滑。令$\zeta_j=\eta_j/S$，则$\zeta_j\in C_c^\infty(V_j)$，$0\leq\zeta_j\leq1$，并且在$\Omega$上$\sum_j\zeta_j=1$，每点附近只有有限项非零。称$\{\zeta_j\}$为从属于覆盖$\{V_j\}$的\term{单位分解}（partition of unity）。

\begin{theorem}[开集内光滑函数的稠密性]
\label{thm:sobolev-meyers-serrin}
设$\Omega\subset\mathbb R^d$为任意开集，$k\geq0$为整数，$1\leq p<\infty$。则$C^\infty(\Omega)\cap W^{k,p}(\Omega)$在$W^{k,p}(\Omega)$中稠密。
\end{theorem}

\begin{proof}
固定$f\in W^{k,p}(\Omega)$与$\theta>0$。

\emph{分块。}令$f_j=\zeta_jf$。由引理~\ref{lem:sobolev-product-rule}，$f_j\in W^{k,p}(\Omega)$，且它和各阶弱导数都在$\operatorname{supp}\zeta_j\subset V_j$之外为零。把$f_j$在$\Omega$外补成零，得到的函数仍记为$f_j$，它属于$W^{k,p}(\mathbb R^d)$：对$\mathbb R^d$上的测试函数$\varphi$，取在$\operatorname{supp}\zeta_j$上为$1$、支撑在$\Omega$内的截断函数$\chi$，则$\int f_j\,\partial^\alpha\varphi=\int f_j\,\partial^\alpha(\chi\varphi)$，而$\chi\varphi$是$\Omega$上的测试函数，所以测试恒等式在全空间成立。

\emph{逐块磨光。}由式\eqref{eq:sobolev-whole-space-mollify}，可取$\varepsilon_j>0$，使$g_j=f_j*\rho_{\varepsilon_j}$满足
\[
\lVert g_j-f_j\rVert_{W^{k,p}(\mathbb R^d)}<\theta\,2^{-j}.
\]
$g_j$的支撑落在$\operatorname{supp}\zeta_j$的$\varepsilon_j$邻域内。$\operatorname{supp}\zeta_j$是$V_j$中的紧集，而$V_j\subset W_j:=G_{j+2}\setminus\overline{G_{j-2}}$且$\overline{V_j}$与$W_j$的补集有正距离，再缩小$\varepsilon_j$，可使$\operatorname{supp}g_j\subset W_j$。误差按$2^{-j}$分配，是为了让总和不超过$\theta$。

\emph{拼合。}令$g=\sum_jg_j$。每点有邻域只与有限多个$W_j$相交，所以$g$在这个邻域上是有限个光滑函数之和，$g\in C^\infty(\Omega)$。同样，在每个$G_m$上，$f=\sum_jf_j$也只是有限和。于是
\[
\lVert g-f\rVert_{W^{k,p}(G_m)}\leq\sum_j\lVert g_j-f_j\rVert_{W^{k,p}(\mathbb R^d)}<\theta .
\]
令$m\to\infty$，由单调收敛定理得$\lVert g-f\rVert_{W^{k,p}(\Omega)}\leq\theta$，特别地$g\in W^{k,p}(\Omega)$。
\end{proof}

这个定理对任意开集成立，不需要边界的任何正则性。但要读清它给出了什么：逼近函数只在$\Omega$内部光滑，在边界附近可以剧烈振荡；它们一般不是紧支撑的，也不一定能连续地延伸到边界。紧支撑光滑函数在$W^{k,p}(\Omega)$中一般并不稠密，第~\ref{sec:sobolev-endpoint-values}小节将看到，$(0,1)$上的常数$1$与所有紧支撑光滑函数的$H^1$距离至少为$1/\sqrt2$。若要求逼近函数光滑到边界，需要先把函数延拓到区域之外，这是下一节的任务。

指数条件$p<\infty$不能去掉。考虑$(0,1)$上的$f(x)=|x-1/2|$与任一$v\in C^1(0,1)\cap W^{1,\infty}(0,1)$。$v'$在$1/2$处连续，所以在$1/2$的足够小邻域内$|v'(x)-v'(1/2)|<\delta$；而$f'$在该点左侧为$-1$、右侧为$1$，因此
\[
\lVert v'-f'\rVert_\infty\geq\max\bigl(|v'(1/2)+1|,\,|v'(1/2)-1|\bigr)-\delta\geq1-\delta .
\]
$\delta$可以任意小，于是$\lVert v-f\rVert_{W^{1,\infty}}\geq1$：本质上确界范数要求导数在每个小邻域都一致接近，连续导数无法一致逼近跳跃的导数。
